\section{Evaluation}
\label{sec:evaluation}

The evaluation is organized around six experiments with fixed agent versions,
pinned model snapshots, and reproducible logs.

\subsection{Research Questions}

The evaluation targets six questions.

\begin{description}
\item[RQ1: Bypass coverage.] Does OS-level enforcement catch the same policy
violation through direct tool calls, shell strings, subprocesses, SDK paths, and
direct syscalls?
\item[RQ2: Corrective feedback.] Does semantic feedback improve task completion
and reduce repeated violations compared with failure without a reason?
\item[RQ3: False positives.] Do benign tasks and sanctioned gate/declassify
paths proceed without unnecessary failure?
\item[RQ4: Overhead.] What are the per-event and end-to-end costs of label
propagation and rule checks?
\item[RQ5: Real-policy coverage.] How many real instructions from the corpus
are observable, expressible, robust, and temptable?
\item[RQ6: Baselines.] Which capabilities are unique relative to tool-layer
agent guardrails and kernel provenance/enforcement systems?
\end{description}

\subsection{Experiments}

\begin{table*}[t]
\centering
\caption{Planned evaluation matrix. Result fields remain \textsc{TBD}.}
\Description{A table listing planned experiments and planned result fields.}
\begin{tabular}{p{0.12\textwidth}p{0.26\textwidth}p{0.26\textwidth}p{0.24\textwidth}}
\toprule
Experiment & Method & Primary metrics & Result status \\
\midrule
E1: Cross-path coverage &
Run the same forbidden action through tool call, nested shell, subprocess, SDK/helper, and direct syscall triggers. Compare against tool-layer baselines. &
Coverage matrix over policy by path; feedback delivery. &
\textsc{TBD} \\
\midrule
E2: Feedback loop &
Run Codex and Claude-style agents under prompt-only, audit, fail-without-reason, and fail-with-feedback conditions. &
Task completion, repeated violations, repair iterations, token/turn cost. &
\textsc{TBD} \\
\midrule
E3: False positives &
Run allowed tasks and sanctioned paths: redaction, test-before-commit, confirmation, and lineage mediation. &
False positive rate, allowed-path success. &
\textsc{TBD} \\
\midrule
E4: Overhead &
Microbench fork/exec/read/write/connect and macrobench real agent tasks with and without ActPlane. &
p50/p99 latency, wall time, memory/map usage, rule-table scaling. &
\textsc{TBD} \\
\midrule
E5: Corpus coverage &
Double-code real instruction-file rules with D1--D7, report agreement, and derive policies/tasks mechanically. &
Observable, expressible, robust, temptable fractions; Cohen's kappa. &
\textsc{TBD} \\
\midrule
E6: Baseline comparison &
Compare against CamQuery, Tetragon, Progent, AgentSpec, and built-in agent sandboxes where runnable. &
Layer, enforce/detect, cross-channel taint, bypass resistance, feedback. &
\textsc{TBD} \\
\bottomrule
\end{tabular}
\label{tab:eval-plan}
\end{table*}

\subsection{Feedback Conditions}

The agent-facing experiment uses four conditions to isolate the contribution of
semantic feedback.

\begin{itemize}
\item C1: prompt-only. The natural-language rule is present, but there is no
ActPlane enforcement.
\item C2: audit. ActPlane observes and reports, but does not fail the action.
\item C3: fail without ActPlane reason. The violating action fails through
\texttt{-EPERM} or \texttt{SIGKILL}, but the agent receives no structured
remediation.
\item C4: fail with feedback. The action fails and ActPlane injects the rule
reason and remediation through the feedback file and hook bridge.
\end{itemize}

The main comparison is C3 to C4: the net benefit of the corrective feedback
loop. C1 to C2 estimates how often the prompt-only harness fails; C2 to C3
separates observation from action failure.

\subsection{Reproducibility}

The artifact includes policies, trigger scripts, agent prompts, raw logs,
statistical notebooks, and a container description pinning kernel, BPF, agent,
and model versions. All reported numbers can be regenerated from these
artifacts.


\section{Related Work}
\label{sec:related}

\subsection{In-Kernel Provenance and Information Flow}

ActPlane builds on a long line of provenance and information-flow systems.
PASS made provenance a storage-layer primitive~\cite{muniswamy2006pass}.
Hi-Fi and LPM used LSM hooks to capture trustworthy system
provenance~\cite{pohly2012hifi,bates2015lpm}. CamFlow captured whole-system
provenance streams over processes, files, and sockets~\cite{pasquier2017camflow}.

CamQuery is the closest prior system~\cite{pasquier2018camquery}. It executes
provenance queries inline in the kernel, propagates labels across process, file,
and network edges, and can prevent policy violations. ActPlane does not claim
that this mechanism is new. The difference is the domain and substrate:
ActPlane targets cooperative AI agents, compiles an agent-oriented source/sink
DSL to eBPF/BPF-LSM, and returns corrective feedback to the actor that violated
the rule.

TaintDroid, Dytan, libdft, and Panorama established dynamic taint tracking at
runtime, binary-instrumentation, or emulation layers~\cite{enck2010taintdroid,
clause2007dytan,kemerlis2012libdft,yin2007panorama}. ActPlane uses a much
coarser object-level taint. That coarseness is what allows kernel-wide coverage,
but it creates an over-taint risk that must be measured.

\subsection{eBPF and LSM Enforcement}

BPF-LSM, originally proposed through KRSI, lets eBPF programs attach to Linux
security hooks and return allow/deny verdicts~\cite{singh2019krsi,linuxbpflsm}.
ActPlane uses this as its pre-operation backend. Contemporary eBPF security
tools and proposals such as Tetragon, Tracee, and eBPF-PATROL demonstrate
practical event capture and enforcement, but generally evaluate per-event
predicates or lineage flags rather than typed cross-channel information
flow~\cite{ciliumtetragon,aquasecuritytracee,ghimire2025ebpfpatrol}.

\subsection{Agent Guardrails}

AgentSpec and Progent are the closest agent-policy systems at the tool
layer~\cite{wang2025agentspec,shi2025progent}. They provide deterministic
checks over tool names, arguments, or framework actions. ActPlane complements
these systems by moving the enforcement point below the tool API. If all
effects pass through the framework, tool-layer policies are sufficient and
cheaper. If the agent can execute arbitrary code, spawn subprocesses, or use
direct system interfaces, the OS layer is the robust boundary.

FIDES and CaMeL show that typed information-flow and capability separation are
useful abstractions for LLM agents~\cite{costa2025fides,debenedetti2025camel}.
They operate inside the agent loop or scaffolding. ActPlane applies the same
style of information-flow reasoning at the OS boundary.

\subsection{Agent Instruction Files}

Recent empirical work studies \texttt{CLAUDE.md}, \texttt{AGENTS.md}, and other
agent context files: their structure, content, maintenance, and effect on task
efficiency~\cite{chatlatanagulchai2025manifests,chatlatanagulchai2025agentreadmes,
santos2025claudeconfig,lulla2026agentsmd}. ActPlane's corpus study asks a
different question: which instructions are behavioral or flow invariants, and
which of those can be observed and enforced below the tool layer?


\section{Conclusion}

ActPlane moves agent harness enforcement below the tool layer. It compiles
agent-declared behavioral contracts into eBPF/BPF-LSM labeled information-flow
rules, propagates labels across process, file, and endpoint objects, and
returns kernel-detected violations as corrective feedback to the agent. Unlike
a classic reference monitor that only denies, ActPlane targets cooperative
agents: it provides audit-only guidance, harness-level termination, and
remediation text that lets the agent complete the original task through an
allowed path. The labeled information-flow mechanism is not new---ActPlane's
contribution is unifying it with an agent-oriented contract language and a
feedback loop on the modern eBPF/BPF-LSM substrate, bridging the
intent--behavior gap that existing single-level harnesses leave open.
